\section{Implementation and evaluation details}
\label{sec:method-contract}
\subsection{Retrieval and repair settings}
The implemented witness-retrieval path selects the top three same-repository
sources using Qwen3-Embedding-8B and FAISS. It returns complete records and
typed input relations. The broader comparison uses the existing 273-source
bank plus validated pre-base additions, with the bank and adapters fixed
before evaluation. Iterative search changes the selection procedure;
sequential updating changes when new experience enters memory. Both use the
same requirement representation and composition interface.

The experiment configuration records source and target IDs, model and agent
versions, index and adapter versions, prompts, budgets, and verifiers.
This makes the retrieval, behavioral, and cost comparisons reproducible.

\subsection{Preparing the current executable program}
\label{sec:public-program}
The development range example extracts model and test snippets from the
public issue and adds researcher-written setup and observation code. The
range and timezone adapters are also researcher-written. \expmark{The
following automatic preparation and applicability workflow is specified for
the additional evaluation and still requires implementation and validation.}

Given an issue and the original repository, the coding agent first inspects
public code, tests, and documentation. It extracts or constructs a reproducer,
prepares dependencies and objects, and identifies the operation to combine.
The output contains setup, exercise, and observation functions; the execution
command and environment; input roles; and the public evidence for the expected
outcome. An initial execution checks that the program reaches the intended
operation. A failure of the reported behavior is useful; a setup failure
means preparation is incomplete.

This preparation runs before memory is supplied. All methods receive the
same program, execution record, or failure status for each task and repetition.
Its full resource use counts against every method's total budget. When the
artifact is physically reused, the campaign ledger records its actual charge
once while the per-method estimate includes the full preparation cost.
Preparation has fixed tool, token, and time limits. It uses public information
and leaves the candidate implementation unchanged. If it fails, the agent
continues ordinary repair work, with any applicable requirement supplied as
text. Researchers do not prepare individual fixtures in this automatic setting.

\subsection{Deciding whether history still applies}
The assessor compares the retrieved requirement with the current issue,
version-matched documentation and tests, and reachable pre-base history.
Acceptance requires compatible input and operation semantics and current
evidence for the expected behavior. An explicit behavior change rejects the
old requirement; missing or conflicting evidence produces an uncertain
decision. Only accepted requirements become executable checks. An accepted
requirement without a supported adapter is delivered as cited text.

Each decision records the supporting source and current evidence, the input
mapping, and the reason. The automatic assessor uses the same backbone as
the solver. Independent annotators evaluate these decisions and the resulting
repairs; their labels and checks remain private. Appendix
\ref{sec:applicability-audit} describes this evaluation.

\subsection{Composition interface}
The composer takes the current program $P_q$, memory $m$, typed bindings,
and an adapter registry. Bindings describe constructors, input values and
types, operation arguments, current objects, and their source locations.
It returns a check, its observer, the mapping used, and a construction status.
The steps are:
\begin{enumerate}
\item Match an adapter to the current operation and its type requirements.
\item Map current objects to historical input roles, preserving both triggers
and the current operation sequence.
\item Construct the program and check its fixture, bindings, and observer.
The expected outcome comes from the requirement, independently of the candidate patch.
\item Execute the check and return a behavioral failure as repair feedback.
Record setup or observation errors separately so they can be distinguished
from failures of the requirement.
\end{enumerate}

The range adapter transforms the public query's syntax tree. It preserves
the lazy User and nested query while changing the terminal lookup to a
point range. Equality and range must select the same normalized non-null
value. The timezone adapter instead constructs combinations of database,
ambient, and destination roles. Source extractors recover the bindings these
adapters support. Adding an operation family requires another adapter and
its development tests.

\subsection{Coverage and mechanism comparisons}
\label{sec:support-strata}
Before any run, classify each target from its public issue and repository
as supported by the adapters, outside their scope, or uncertain. Retain the
same groups when reporting preparation, delivery, repair success, and cost.
Failures and text fallback stay in their assigned groups.

The relation ablation gives the same composer either stored bindings or
bindings reconstructed from text. Both use the same record schema, source
IDs, applicability decisions, validator, and budget. Withheld relations are
also omitted from tool explanations. Another contrast supplies the same
bindings to an adapter or a free-form composer. A third supplies execution
observations before editing; ordinary execution tools remain available to
all arms. The separate-versus-joint comparison holds two check slots and
feedback rounds fixed. These comparisons measure the role of each choice
on the development tasks (Table~\ref{tab:planned-graph}).
